\clearpage
\section{符号说明}
\begin{table}[!htbp]\centering
\caption{主要符号及单位}\label{tab:symbols}
\begin{tabular}{cll}\toprule
符号&含义&单位\\\midrule
$r,z$&径向、轴向坐标；$z=0$ 为中截面&m\\
$R(t),L$&当前半径、长度&m\\
$T,T_e$&药材、环境温度&$^\circ$C\\
$C$&药材干基含水率&kg水/kg干物质\\
$C_e$&题给环境有效含湿驱动值&kg/kg\\
$\rho,c_p,k$&经验密度、比热容、导热系数&kg/m$^3$，J/(kg K)，W/(m K)\\
$\rho_d$&单位当前体积的干物质质量&kg干物质/m$^3$\\
$D$&有效水分扩散系数&m$^2$/s\\
$h_T,h_m$&表面换热、有效传质系数&W/(m$^2$K)，m/s\\
$\xi$&材料坐标 $r/R(t)$&1\\
$v_r,w_r$&材料、网格径向速度&m/s\\
$V_i,A_f$&单位长度控制体体积、界面面积&m$^2$，m\\
$t_d$&最大含水率达到阈值的首达估计&s\\\bottomrule
\end{tabular}\end{table}
扩散系数中的绝对温度为 $T_K=T+273.15$；温差的摄氏度与开尔文数值相同。本文规定时间首达误差与场量误差分别报告，不以提交表的四位小数格式代替精度说明。
